
\documentclass{article}
\usepackage[left=2cm, right=2cm]{geometry}
\usepackage{amsmath}

\title{Assignment 6}
\author{Christian Parker \\ CXP22004}

\begin{document}

\maketitle

\begin{enumerate}
    % Problem 1
    \item
    \begin{enumerate}
        \item
        \begin{align*}
            J(u) &= \int_{0}^{1} (u(x) - u_0 (x))^2 dx + \lambda \int_{0}^{1} (u'(x))^2 dx \\
            &= \int_{0}^{1} (u - u_0)^2 + \lambda (u')^2 dx \\
            F(x, u, u') &= (u - u_0)^2 + \lambda (u')^2 \\
            \frac{\delta F}{\delta u} &= 2(u - u_0) \\
            \frac{\delta F}{\delta u'} &= 2 \lambda u'
        \end{align*}

        The Euler-Lagrange equation is:
        \begin{align*}
            \frac{\delta F}{\delta u} - \frac{d}{dx} \frac{\delta F}{\delta u'} &= 0 \\
            2(u - u_0) - \frac{d}{dx} 2 \lambda u' &= 0 \\
            u - u_0 - \lambda u'' &= 0 \\
            u - \lambda u'' &= u_0
        \end{align*}
    \end{enumerate}

    % Problem 2
    \item
    \begin{align*}
        E(u) &= \alpha \int_{\Omega} | \nabla u | dxdy + \int_{\Omega} (u - u_0 \cdot ln(u)) dxdy \\
        &= \int_{\Omega} \alpha | \nabla u | + u - u_0 \cdot ln(u)\ dxdy \\
        F(u, u_x, u_y) &= \alpha | \nabla u | + u - u_0 \cdot ln(u) \\
        \frac{\delta F}{\delta u} &= 1 - \frac{u_0}{u} \\
        \frac{\delta F}{\delta \nabla u} &= \left( \frac{\delta F}{\delta u_x}, \frac{\delta F}{\delta u_y} \right) \\
        &= \left( \alpha \frac{\delta}{\delta u_x} \sqrt{u_x^2 + u_y^2}, \alpha \frac{\delta}{\delta u_y} \sqrt{u_x^2 + u_y^2} \right) \\
        &= \alpha \left( \frac{2 u_x}{2 \sqrt{u_x^2 + u_y^2}}, \frac{2 u_y}{2 \sqrt{u_x^2 + u_y^2}} \right) \\
        &= \alpha \frac{(u_x, u_y)}{\sqrt{u_x^2 + u_y^2}} \\
        &= \alpha \frac{\nabla u}{| \nabla u |}
    \end{align*}

    The Euler-Lagrange equation is:
    

    \begin{align*}
        \frac{\delta F}{\delta u} - \operatorname{div} \left( \frac{\delta F}{\delta \nabla u} \right) &= 0 \\
        1 - \frac{u_0}{u} - \operatorname{div} \left( \alpha \frac{\nabla u}{| \nabla u |} \right) &= 0 \\
    \end{align*}

    % Problem 3
    \item
    \begin{enumerate}
        \item[(c)]
        Graph cut performed the best, providing a clear outline of the tumor region. The Image
        Segmenter required drawing lines on the foreground and background before running the
        graph cut algorithm. The drawing I made outlined the tumor fairly well. This headstart
        is likely why this algorithm performed so well.

        Both active contour methods performed worse than the graph cut method.
        The Chan-Vese method was highly sensitive to the input mask.
        If the mask was too big, it would mask the entire brain rather than just the tumor region.
        In contrast, the edge method mask remained within the tumor for every mask I tried.
        However, the mask was always noticeably smaller than the tumor.
        
        With a proper mask, the Chan-Vese method provided better coverage than the edge method.
        The tumor was generally outlined well, with some additional tendrils
        reaching out to nearby bright spots that weren't part of the tumor.
    \end{enumerate}

\end{enumerate}

\end{document}